% Ported verbatim from paper/top3_joint_moment_bound.tex lines 177-206.
% Only cross-reference wording was adapted; no mathematics was altered.
The proof has two parts. First, second variation on the actual probability
space gives a negative semidefinite kernel. On each of four sign classes,
this controls the rank of the positive-part Gram kernel and forces finite
canonical support. Moment-preserving variations then imply vector
relations against every binary row of the maximizing kernel. Those rows
separate distinct support points. A minimal-dependency argument shows
that no second-moment relation can survive on a positive clique.

The variation, cloning, and orthant-reduction mechanisms have close
antecedents in Kumar, Mohar, Mojallal and Pragada~\cite{KMMP}. Related graphon
ellipse and clique-cloning arguments for the two-eigenvalue problem
appear in Kumar, Liu, Monterde, Pragada and Tait~\cite{KLMPT}. General
bounds for minimal dependencies of powers of linear forms are due to
Li and Liu~\cite{LiLiu}. Our separating-observation lemma is stated and
proved with its specific hypotheses in Section~\ref{sec-separation};
we do not claim a new general bound for such dependencies.

Second, for one isotropic law we retain the mean deficit, radial variance,
regression error, and weighted harmonic moments of degrees three and four.
A signed quartic majorant couples them in one exact identity. Directional
Hermite inequalities and a radial comparison lead to an explicit sum of
nonnegative terms. Harmonic moment and polynomial-majorant methods are
classical; see Glazyrin~\cite{Glazyrin}, Der\k{e}gowska and
Lewandowska~\cite{DL,DL-SW}, and Sivashankar, Tang and Wakhare~\cite{STW}.
The contribution here is the particular joint control of these errors
for the mixed signed and absolute energy, rather than the use of
polynomial majorants by itself. The support theorem supplies a small
additional term through self-pairs of actual atoms. Diagonal removal and
effective-cardinality estimates are standard weighted-energy devices;
see~\cite{Weighted}.
